\begin{figure*}[t]
  \centering
  \resizebox{\textwidth}{!}{%
  \begin{tikzpicture}[
    x=0.8cm,y=1.0cm,
    block/.style={draw, rounded corners=2pt, align=center, minimum height=10mm,
      text width=32mm, fill=blue!4},
    flow/.style={-{Latex[length=2.2mm]}, thick}
  ]
    \begin{scope}[shift={(4.6,0)}]
      \draw[->] (0,0) -- (7.1,0) node[right] {层 $l$};
      \draw[->] (0,0) -- (0,2.5) node[above] {$\lambda_{c,l}$};
      \draw[dashed, gray] (1.3,0) -- (1.3,0.55);
      \draw[dashed, gray] (3.3,0) -- (3.3,2.05);
      \draw[dashed, gray] (5.3,0) -- (5.3,0.55);
      \draw[very thick, blue, dashed] (0,0) -- (1.3,0);
      \draw[very thick, blue, dashed] (1.3,0.5) -- (3.3,2.0) -- (5.3,0.5);
      \draw[very thick, blue, dashed] (5.3,0) -- (6.7,0);
      \draw[blue, fill=white, very thick] (3.3,2.0) circle (1.8pt);
      \filldraw[blue] (1.3,0.5) circle (1.6pt);
      \filldraw[blue] (5.3,0.5) circle (1.6pt);
      \draw[blue, fill=white, very thick] (1.3,0) circle (1.6pt);
      \draw[blue, fill=white, very thick] (5.3,0) circle (1.6pt);
      \foreach \x/\y in {1/0,2/1.025,3/1.775,4/1.475,5/0.725,6/0}
        \filldraw[black] (\x,\y) circle (1.5pt);
      \node[below] at (1.3,0) {$p_c-d_c$};
      \node[below] at (3.3,0) {$p_c$};
      \node[below] at (5.3,0) {$p_c+d_c$};
      \node[left] at (0,2.0) {$M_c$};
      \node[left] at (0,0.5) {$m_c$};
      \node[align=center, font=\scriptsize] at (3.4,2.65)
        {虚线为连续核；黑点为整数层实际采样\\
        $p_c\notin\mathbb{Z}$ 时实际层不取得 $M_c$};
      \node[align=left, font=\scriptsize] at (7.7,1.1)
        {$m_c>0$ 时边界处跳变\\各组件独立采样一组参数};
    \end{scope}

    \begin{scope}[shift={(0,-4.3)}]
      \node[block] (directions) {相邻层有效方向\\
        $q_{\lfloor j\rfloor},q_{\lceil j\rceil}$};
      \node[block, right=8mm of directions] (interp) {线性插值后归一化\\
        $v=\operatorname{norm}($\\$\operatorname{lerp}(\cdot))$};
      \node[block, right=8mm of interp] (project) {秩一投影项\\
        $v(v^{\top}W_{c,l})$};
      \node[block, right=8mm of project] (updated) {组件与层特定更新\\
        $W'_{c,l}=W_{c,l}-$\\$\lambda_{c,l}v(v^{\top}W_{c,l})$};
      \draw[flow] (directions) -- node[above, font=\scriptsize]
        {浮点 direction index} (interp);
      \draw[flow] (interp) -- (project);
      \draw[flow] (project) -- node[above, font=\scriptsize]
        {$\lambda_{c,l}$} (updated);
      \node[draw, dashed, fit=(project)(updated), inner sep=4mm,
        label={[font=\scriptsize]below:FULL 模式随后恢复行范数并以随机低秩分解近似差值}] {};
    \end{scope}
  \end{tikzpicture}%
  }
  \caption{上：组件 $c$ 的连续层权重核及其整数层采样；不同组件独立采样核参数。
  下：浮点方向索引、矩阵投影与层权重的关系。FULL 行归一化在秩一更新之后
  引入非线性，因此最终 LoRA 只近似完整差值。}
  \label{fig:weight-kernel}
\end{figure*}
